\section{Related Work}
With relatively recent popularity of fairness in machine learning and natural language processing domains, the need to find a universal and a more complete fairness definition and measure is crucial. Although finding such definition and measure is a challenge not only in machine learning but also in social and political sciences, steps need to be taken to make current definitions evolve and cover more real world cases. In light of this many fairness definitions have been proposed. Some tried to complement others and some starting a new direction and view-point on their own. Different body of work tried to incorporate the proposed definitions in different downstream tasks such as classification and regression \cite{pmlr-v81-menon18a,
berk2017convex, Krasanakis:2018:ASR:3178876.3186133, agarwal2019fair,goel2018non}.
  
\subsection{Fairness Definitions}
For a more complete list of existing fairness definitions there exists papers that survey \cite{mehrabi2019survey} and explain \cite{verma2018fairness} proposed definitions. Here we will elaborate some important and widely known definitions related to our work introduced in this paper.

\subsubsection{Statistical Parity} In statistical parity \cite{Dwork:2012:FTA:2090236.2090255} the goal is to satisfy $P(\hat{Y} |A = a) = P(\hat{Y}|A = b)$. This notion states that regardless of the belonging demographic group $A$, the predicted outcome should be the same for both demographic groups $A=a$ and $A=b$.
\subsubsection{Equalized Odds} In equalized odds \cite{hardt2016equality}   the goal is to satisfy  $P(\hat{Y}=1|A=a,Y=y) = P(\hat{Y}=1|A=b,Y=y)$ for $y\in\{0,1\}$. This notion states that both groups $A=a$ and $A=b$ should have equal true positive and false positive rates.
\subsubsection{Equal Opportunity} In equal opportunity \cite{hardt2016equality} the goal is to satisfy $P(\hat{Y}=1|A=a,Y=1) = P(\hat{Y}=1|A=b,Y=1)$. This notion states that both demographic groups $A=a$ and $A=b$ should have equal true positive rates.
\subsubsection{Counterfactual Fairness} In counterfactual fairness \cite{NIPS2017_6995} the goal is to satisfy
$P(\hat{Y}_{A\xleftarrow{}a }(U)=y|X=x,A=a) = P(\hat{Y}_{A\xleftarrow{}a'}(U)=y|X =x,A=a)$ under any $X=x$ and $A=a$ for all $y$ and for any value $a'$ feasible for $A$. The perception in counterfactual fairness definition is that if the decision is the same in both the actual and counterfactual world where the individual belonged to a different group then it is called a fair decision. 
% \subsubsection{Calibration???? Do I want to talk about this of fairness through unawareness????} \cite{chouldechova2017fair}

\subsection{Fairness Definitions in Classification}
Research in fairness domain does not conclude itself in defining fairness definitions and measures but also incorporation of these definitions in tasks such as classification \cite{calders2010three,huang2019stable}. This incorporation on its own is a challenge. Some methods introduce pre-processing methods that augment the train data for discrimination removal \cite{Kamiran2012}. Some other methods perform in-processing which tries to incorporate the fairness objective during training phase \cite{kamishima2012fairness,zafar2015fairness,wu2018fairnessaware}. Other definitions require post-processing techniques in which discrimination removal is performed after the training phase treating the model as a black box~\cite{hardt2016equality,NIPS2017_7151}. Literature on fair classification also targets a wide variety of fairness definitions, such as equality of opportunity and equalized odds \cite{hardt2016equality,woodworth2017learning}, statistical parity \cite{agarwal2018reductions}, and subgroup fairness \cite{pmlr-v97-ustun19a}, and their incorporation in classification task.